\documentclass[10pt]{mypackage}

\usepackage[uprightgreeks,light]{kpfonts}
\renewcommand{\coloneq}{\coloneqq}
%\usepackage{homework}
\usepackage{notes}

\usepackage[ backend=bibtex, style = alphabetic, sorting=ynt ]{biblatex}
\addbibresource{all_references.bib}

\usepackage{parskip}

\fancyhf{}
\fancyhead[R]{Avinash Iyer}
\fancyhead[L]{Central Limit Theorem(s)}
\fancyfoot[C]{\thepage}

\setcounter{secnumdepth}{0}

\begin{document}
\RaggedRight
The central limit theorem is one of the most celebrated theorems in probability. Here, we will discuss the central limit theorem both in classical probability and in free probability.
\section{Classical Central Limit Theorem}%
The classical central limit theorem is as follows.
\begin{theorem}[Classical Central Limit Theorem]
  Let $\left( X_n \right)_{n\in \N}$ be a sequence of independent and identically distributed random variables with common mean $\mu$ and variance $\sigma^2$. Then,
  \begin{align*}
    \frac{\left( X_1 + \cdots + X_N \right) - N\mu}{\sigma\sqrt{N}} &\rightarrow \mathcal{N}\left( 0,1 \right)
  \end{align*}
  in distribution, where $\mathcal{N}\left( 0,1 \right)$ has distribution function given by
  \begin{align*}
    d\mu &= \frac{1}{\sqrt{2\pi}} e^{-t^2/2}\:dt.
  \end{align*}
\end{theorem}
\subsection{Overview of Probability Theory}%

\subsection{Proof via Fourier Analysis}%
\subsection{Proof via Moments and Partitions}%

\section{Free Central Limit Theorem}%
\subsection{Overview of Noncommutative Probability}%
\subsection{Proof using Moments and Noncrossing Partitions}%
\subsection{Free Cumulants and the $R$-Transform}%

\end{document}
